\chapter{Теоретические сведения}
\label{ch:theory}

\section{Figma как инструмент UI/UX-дизайна}

\textbf{Figma} — облачный векторный графический редактор интерфейсов,
разработанный компанией \textit{Figma, Inc.} в 2016~году. Ключевые особенности:
\begin{itemize}
  \item \textbf{Облачное хранение} — файлы лежат на серверах Figma,
        доступны через браузер с любого устройства;
  \item \textbf{Совместная работа в реальном времени} — несколько
        дизайнеров могут одновременно редактировать один файл,
        видя курсоры друг друга;
  \item \textbf{История версий} — каждое изменение сохраняется, есть
        возможность отката к любой ранней редакции;
  \item \textbf{Плагины} — расширяют возможности (поиск изображений,
        генерация цветовых палитр, экспорт в код и т.\,п.);
  \item \textbf{Кроссплатформенность} — веб-версия работает в любом
        современном браузере, есть также desktop-приложения для
        Windows и macOS.
\end{itemize}

\section{Структура файла Figma}

Файл Figma организован иерархически:
\begin{itemize}
  \item \textbf{File} — отдельный документ (\texttt{*.fig}), содержит одну
        или несколько страниц;
  \item \textbf{Page} (\texttt{Страница}) — независимый холст; обычно одна
        страница соответствует одному смысловому разделу проекта
        (например, главная страница, корзина, профиль);
  \item \textbf{Frame} (\texttt{Фрейм}) — прямоугольная область, которая
        одновременно является контейнером для дочерних слоёв и единицей
        экспорта. Фрейм похож на artboard в Adobe Illustrator или
        \texttt{<div>} в HTML;
  \item \textbf{Layer} (\texttt{Слой}) — отдельный графический объект:
        прямоугольник, текст, изображение, эллипс и т.\,п.;
  \item \textbf{Group} (\texttt{Группа}) — объединение нескольких слоёв
        без визуального контейнера; в отличие от фрейма группа не имеет
        собственных размеров и заливки.
\end{itemize}

\textbf{Разница фрейм/группа.} Фрейм — самостоятельный элемент с заданной
шириной, высотой и стилями; в нём работают \texttt{Constraints} и
\texttt{Layout grid}. Группа — лишь логическое объединение слоёв,
изменяющее размер вместе с содержимым. В макетах веб-страниц
рекомендуется использовать именно фреймы.

\section{Геометрические объекты и работа с цветом}

Базовые геометрические инструменты: \texttt{Rectangle}~(\texttt{R}),
\texttt{Ellipse}~(\texttt{O}), \texttt{Line}~(\texttt{L}),
\texttt{Polygon}, \texttt{Star}. Параметры объекта расположены на правой
панели:
\begin{itemize}
  \item \textbf{Fill} — заливка: одноцветная, градиентная (linear,
        radial, angular, diamond) либо изображение;
  \item \textbf{Stroke} — обводка: цвет, толщина, положение
        (\textit{inside}, \textit{outside}, \textit{center}),
        пунктир/штрих;
  \item \textbf{Corner radius} — скругление углов прямоугольников
        (общее или для каждого угла отдельно);
  \item \textbf{Opacity} — прозрачность объекта (\texttt{0--100\,\%});
  \item \textbf{Effects} — тени (\textit{drop/inner shadow}) и размытие
        (\textit{layer/background blur}).
\end{itemize}
Цвета задаются в HEX, RGB или HSL; библиотека стилей позволяет сохранять
часто используемые цвета как переменные.

\section{Текстовые слои}

Инструмент \texttt{Text}~(\texttt{T}) создаёт текстовый слой. Свойства:
\begin{itemize}
  \item шрифт и начертание (Figma подтягивает Google~Fonts и системные
        шрифты);
  \item размер кегля, межбуквенное расстояние, высота строки;
  \item выравнивание по горизонтали и вертикали;
  \item поведение блока: \textit{Auto~width}, \textit{Auto~height},
        \textit{Fixed~size}.
\end{itemize}
Текстовые стили (\texttt{Text styles}) сохраняются в библиотеку файла
и применяются ко множеству слоёв одним кликом.

\section{Работа с изображениями. Плагин Unsplash}

Растровое изображение в Figma — это \texttt{Fill} типа \texttt{Image}
для любой геометрической фигуры. Импорт: \texttt{Place~image}
(\texttt{Ctrl+Shift+K}) или перетаскивание файла на холст.

\textbf{Unsplash} — плагин, открывающий каталог бесплатных фотографий
сервиса \texttt{unsplash.com} прямо в Figma. Установка: меню
\texttt{Plugins~\(\rightarrow\)~Browse plugins in Community~\(\rightarrow\)
Unsplash~\(\rightarrow\)~Install}. Использование: выделить целевой
объект, запустить плагин, выбрать или найти фотографию~--- она
автоматически подставляется в \texttt{Fill}.

\section{Constraints (привязки)}

\textbf{Constraints} управляют поведением слоя при изменении размеров
родительского фрейма. Привязки задаются независимо по горизонтали и
вертикали:
\begin{itemize}
  \item \texttt{Left} / \texttt{Right} / \texttt{Top} / \texttt{Bottom} ---
        слой <<прибит>> к соответствующей стороне; расстояние до неё
        сохраняется;
  \item \texttt{Center} — слой остаётся по центру оси при изменении
        ширины/высоты родителя;
  \item \texttt{Scale} — размеры слоя пропорционально подстраиваются;
  \item \texttt{Left and Right} / \texttt{Top and Bottom} — слой
        растягивается, сохраняя оба отступа.
\end{itemize}
Грамотно настроенные привязки позволяют создавать одну версию макета
и быстро адаптировать её под другие ширины фрейма.

\section{Layout grid (сетка)}

\textbf{Layout grid} — направляющая сетка для выравнивания элементов.
Применяется к фрейму; видна на холсте, но не экспортируется. Типы:
\begin{itemize}
  \item \texttt{Grid} — равномерная решётка квадратами; удобна для
        иконок и пиктограмм;
  \item \texttt{Columns} — вертикальные колонки; основной тип для
        веб-макетов;
  \item \texttt{Rows} — горизонтальные строки.
\end{itemize}
Параметры колоночной сетки: \texttt{Count} (число колонок),
\texttt{Width} (ширина колонки или \texttt{Auto}), \texttt{Margin}
(внешний отступ), \texttt{Gutter} (промежуток между колонками).
Стандарт веб-дизайна — 12-колоночная сетка с отступами и гаттером
\(16\text{--}32\,\)px.

\section{Адаптивный дизайн}

\textbf{Адаптивный дизайн} — подход, при котором интерфейс корректно
отображается на устройствах разной ширины. Ключевые брейкпоинты:
\begin{itemize}
  \item \texttt{Desktop} — от \(1280\,\)px (типично \(1440\,\)px);
  \item \texttt{Tablet} — от \(768\,\)px до \(1279\,\)px
        (типично \(768\,\)px);
  \item \texttt{Mobile} — до \(767\,\)px (типично \(360\,\)px).
\end{itemize}
Приёмы адаптации при сужении фрейма:
\begin{itemize}
  \item уменьшение размеров шрифтов и отступов;
  \item перестроение многоколоночной сетки в одно- или двухколоночную;
  \item скрытие второстепенных элементов (например, развёрнутого меню
        в пользу <<гамбургер>>-иконки);
  \item замена крупных декоративных изображений на компактные.
\end{itemize}
В Figma адаптация выполняется копированием макета и изменением ширины
корневого фрейма; правильные \texttt{Constraints} автоматически
пересобирают часть контента, остальное корректируется вручную.

\endinput
